\documentclass[a4paper,USenglish,cleveref,autoref,thm-restate]{lipics-v2021}

\bibliographystyle{plainurl}

\title{Artifact: Ownership Refinement Types for Pointer Arithmetic and Nested Arrays}

%% TODO: Update author list to match paper
\author{Yusuke Fujiwara}{Kyoto University}{hujiwara.yuusuke.28n@st.kyoto-u.ac.jp}{}{}
\author{Yusuke Matsushita}{Kyoto University}{ymat@fos.kuis.kyoto-u.ac.jp}{https://orcid.org/0000-0002-5208-3106}{}
\author{Kohei Suenaga}{Kyoto University}{ksuenaga@fos.kuis.kyoto-u.ac.jp}{https://orcid.org/0000-0002-7466-8789}{}
\author{Atsushi Igarashi}{Kyoto University}{igarashi@fos.kuis.kyoto-u.ac.jp}{https://orcid.org/0000-0002-5143-9764}{}


\authorrunning{Y. Fujiwara, Y. Matsushita, K. Suenaga, and A. Igarashi}

\Copyright{Yusuke Fujiwara, Yusuke Matsushita, Kohei Suenaga, and Atsushi Igarashi}

% TODO: Update CCS concepts to match paper
\ccsdesc[500]{Software and its engineering~Automated static analysis}
\ccsdesc[300]{Theory of computation~Type theory}

\keywords{ownership types, refinement types, pointer arithmetic, nested arrays, verification}

%% TODO: Update related article DOI
\RelatedArticle{Yusuke Fujiwara, Yusuke Matsushita, Kohei Suenaga, and Atsushi Igarashi, ``Ownership Refinement Types for Pointer Arithmetic and Nested Arrays'', in Proceedings of ECOOP 2026, LIPIcs, Vol.~XXX, pp.~YY:1--YY:ZZ, 2026.\newline \url{https://doi.org/10.4230/LIPIcs.ECOOP.2026.XXX}}

%% TODO: Update artifact DOI from Zenodo
\RelatedArticleResource{}{Software Heritage \url{https://archive.softwareheritage.org/swh:1:dir:XXXXXXXXXXXXXXXXXXXXXXXXXXXXXXXXXXXXXXXX}}

\begin{document}

\maketitle

\begin{abstract}
This artifact accompanies the paper ``Ownership Refinement Types for Pointer Arithmetic and Nested Arrays.''
It provides \texttt{nested\_array\_ConSORT}, a tool for automated ownership type inference and
refinement type checking for imperative programs with pointer arithmetic and nested arrays.
The artifact includes the tool implementation, benchmark programs, evaluation scripts to
reproduce all experimental results (Tables~1--3) from the paper, and a comparison baseline
(Extended\_ConSORT by Tanaka et al.).
The artifact is packaged as a Docker image with all dependencies pre-installed.
\end{abstract}

% SCOPE

\section{Scope}

The artifact supports the experimental evaluation presented in the paper.
Specifically, it enables reproduction of:

\begin{itemize}
  \item \textbf{Table~1}: Verification timing and alias statement counts for 19 nested array benchmarks.
  \item \textbf{Table~2}: Inferred ownership term solutions (SMT2 output files for manual inspection).
  \item \textbf{Table~3}: Performance comparison with the tool by Tanaka et al.\ on 8 integer array benchmarks.
\end{itemize}

The tool accepts programs written in a custom \texttt{.imp} language and automatically
performs ownership type inference (using Z3) followed by refinement type checking (using hoice).
Both positive examples (expected to verify) and negative examples (expected to fail) are included.

% CONTENT

\section{Content}

The artifact consists of:

\begin{itemize}
  \item A pre-built Docker image (\texttt{nested-array-consort:ecoop26}) containing:
    \begin{itemize}
      \item The \texttt{nested\_array\_ConSORT} tool (OCaml, built with dune)
      \item Z3 SMT solver versions 4.11.2 and 4.14.1
      \item hoice CHC solver v1.10.0
      \item Extended\_ConSORT (Tanaka et al.) for baseline comparison
      \item All benchmark \texttt{.imp} programs
      \item Evaluation scripts for reproducing Tables~1--3
    \end{itemize}
  \item Source code of the tool
  \item This artifact description document
\end{itemize}

% GETTING STARTED

\section{Getting Started Guide}

\subsection{Loading the Docker Image}

\begin{verbatim}
docker load -i nested-array-consort-ecoop26.tar.gz
docker run -it nested-array-consort:ecoop26
\end{verbatim}

This drops you into a shell at \texttt{/home/opam/app/myproject/}.

\subsection{Quick Verification (10 minutes)}

Run the kick-the-tires script:

\begin{verbatim}
bash eval/run_short.sh
\end{verbatim}

This runs a subset of benchmarks from Tables~1 and 3 and prints a pass/fail summary.
All benchmarks should report \texttt{[PASS]}.

\subsection{Running a Single Benchmark}

To verify a single program manually:

\begin{verbatim}
dune exec myproject -- ./example/nested_arrays/swap.imp
\end{verbatim}

Expected output includes \texttt{ownership: sat} and \texttt{refinement: sat}.

% DETAILED EVALUATION

\section{Evaluation Instructions}

\subsection{Full Evaluation}

To reproduce all tables from the paper:

\begin{verbatim}
bash eval/run_all.sh
\end{verbatim}

This takes approximately 30 minutes.
Results are saved as Markdown and CSV files in \texttt{eval/results/}.

\subsection{Individual Tables}

Each table can be reproduced independently:

\begin{center}
\begin{tabular}{llll}
\hline
\textbf{Paper} & \textbf{Script} & \textbf{Output} & \textbf{Time} \\
\hline
Table~1 & \texttt{eval/table1.sh} & \texttt{eval/results/table1.md} & $\sim$10 min \\
Table~2 & \texttt{eval/table2.sh} & \texttt{eval/results/table2/} & $\sim$5 min \\
Table~3 & \texttt{eval/table3.sh} & \texttt{eval/results/table3.md} & $\sim$10 min \\
\hline
\end{tabular}
\end{center}

\subsection{Expected Variation}

\begin{itemize}
  \item \textbf{Absolute timing} will differ from the paper due to hardware differences.
  \item \textbf{Sat/unsat results} must match the paper exactly.
  \item \textbf{Ownership solutions} (Table~2) are deterministic and should match the paper.
\end{itemize}

\section{Notes}

The tool invokes Z3 and hoice as external processes. The evaluation scripts switch between
Z3 versions by manipulating the \texttt{PATH} environment variable. All intermediate solver
files are written to \texttt{experiment/} within the project directory. Due to this shared
state, benchmarks must be run sequentially (not in parallel).

% REFERENCES (optional)
%\bibliography{references}

\end{document}
